\section{相対論的理論への動機}
上で見た無限深井戸での離散スペクトル、有限障壁でのトンネル効果、調和振動子の零点エネルギーのように、非相対論的量子力学は多くの現象を具体的に説明する。しかし、この成功はそのまま相対論的領域へは拡張できない。その理由は二つある。第一に、その基本方程式
\[
i\hbar \frac{\partial}{\partial t}\psi = \hat{H}\psi
\]
は時間に関して1階、空間に関しては一般に2階であり、時間と空間を非対称に扱っている。この形式はガリレイ変換の下では自然だが、ローレンツ対称性とは両立しない。

第二に、ここまでの記述では粒子数を固定したヒルベルト空間を暗黙に前提していた。しかし、相対論と量子論を同時に考えると、十分高いエネルギーでは粒子の生成・消滅が起こりうるため、この前提自体が揺らぐ。

したがって、ここでの最初の課題は、ローレンツ対称性と整合する一粒子波動方程式を作ることである。以降ではまずその試みを行い、その後でその試み自体が場の量子論を要請することを見る。
